\newcommand{\AlignAxis}{\bm{\alpha}}

\subsection{Summary}
Transformation:
\begin{enumerate}
    \item \emph{Update} Call \ref{EuclidIsom::update} to update \(\bm{c},\bm{R}\)
\end{enumerate}

Isometry:
\begin{enumerate}
    \item \emph{Initialize}\label{EuclidIsom::initialize} Compute \(\bm{R}_0\) and \(\bar{\bm{C}}\) from \eqref{eq:init-Cbar}
    \item \emph{Update}\label{EuclidIsom::update} \(\bm{R}\) from nodal rotations \(\bm{R}_a\) and new chord \(\Delta\bm{x}\), then compute \(\bar{\bm{c}}\) using \eqref{eq:solve-cbar}.
\end{enumerate}

\subsection{Local System}
\begin{equation}
\begin{cases}
\mathbf{e}_1=\dfrac{\Delta \bm{X}}{\|\Delta \bm{X}\|} \\
\mathbf{e}_2=\dfrac{\mathbf{v} \times \frac{\Delta \bm{X}}{\|\Delta \bm{X}\|}}{\left\|\mathbf{v} \times \frac{\Delta \bm{X}}{\|\Delta \bm{X}\|}\right\|} \\
\mathbf{e}_3 = \mathbf{e}_1 \times \mathbf{e}_2
\end{cases}
\end{equation}


\subsection{Linear System (A1)}

\subsubsection{Transformation}

Consider a line segment in 3D defined by two endpoints \(\bm{X}_i\) and \(\bm{X}_j\). 
After deformation, the endpoints have moved to \(\bm{X}_i + \bm{u}_i\)
 and \(\bm{X}_j + \bm{u}_j\). 
Let 
\[
\Delta \bm{X} \;\triangleq\; \bm{X}_j - \bm{X}_i
\] 
be the direction of the line in the reference configuration, and 
\begin{equation}
\Delta \bm{x} \;\triangleq\; (\bm{X}_J + \bm{u}_J)\;-\;(\bm{X}_I + \bm{u}_I)
\end{equation}
is the direction in the deformed configuration. 
The line's length originally is 
\[
L_0 \; \triangleq \; \|\Delta \bm{X}\|
\] 
and after deformation the length is 
\[
L_n \; \triangleq \; \|\Delta \bm{x}\|
\]

\paragraph{Translation}
The line is represented in the ``natural'' reference configuration \(\bar{\bm{X}}\) by a Euclidean transformation of the reference placement \(\bm{X}\):
\[
\bar{\bm{X}} = \bm{R}_0^{\mathrm{t}}\bm{X} - \bar{\bm{C}}
\]
Similarly, the deformed configuration may be transformed:
\[
\bar{\bm{x}}(\xi) \triangleq  \bm{R}^{\mathrm{t}}\bm{x} - \bar{\bm{c}}
\]
The deformation from \(\bar{\bm{X}}\) to \(\bar{\bm{x}}\) is characterized by the displacement \(\bar{\bm{u}}\):
\[
\begin{aligned}
\bar{\bm{u}}(\xi) &:=  \bar{\bm{x}} - \bar{\bm{X}} \\
&= \bm{R}^{\mathrm{t}}(\bm{X}(\xi)+\bm{u}(\xi)) - \bar{\bm{c}} - (\bm{R}_0^{\mathrm{t}}\bm{X}(\xi) - \bar{\bm{C}}) \\
&= \bm{R}^{\mathrm{t}}(\bm{X}(\xi)+\bm{u}(\xi)) - \Delta \bar{\bm{c}} - \bm{R}_0^{\mathrm{t}}\bm{X}(\xi) \\
&= \bm{R}^{\mathrm{t}}(\bm{u}(\xi) - \Delta \bm{c}) + \Delta \bm{R}^{\mathrm{t}}\bm{X}(\xi) \\
\end{aligned}
\]
with \(\Delta \bar{\bm{c}} \triangleq \bar{\bm{c}} - \bar{\bm{C}}\) and \(\Delta \bm{R} \triangleq \bm{R} - \bm{R}_0\)
so that
\begin{equation}
    \bm{u}(\xi)\;\mapsto\; \bm{R}^{\mathrm{t}}\bm{u}(\xi) - \Delta \bar{\bm{c}} + \Delta \bm{R}^{\mathrm{t}}\bm{X}(\xi)
\label{eq:iso-u}
\end{equation}
The translations \(\bar{\bm{C}}\) and \(\bar{\bm{c}}\) are selected to enforce \(\bar{\bm{u}}_I = \bar{\bm{x}}_I = \bm{0}\), furnishing for \(\bar{\bm{c}}\)
%
\begin{equation}
\bar{\bm{x}}_I = \bm{0} = \bm{R}^{\mathrm{t}}\bm{X}_I + \bm{R}^{\mathrm{t}}(\bm{u}_I - \bm{c})
\quad\implies\quad 
\bar{\bm{c}} = \bm{R}^{\mathrm{t}}(\bm{X}_I + \bm{u}_I)
\label{eq:solve-cbar}
\end{equation}
%
To find \(\bar{\bm{C}}\) use \(\bar{\bm{c}}\) above in \eqref{eq:iso-u}:
\[
\begin{aligned}
\bar{\bm{u}}(\xi)
&= \bm{R}^{\mathrm{t}}(\bm{X}(\xi)+\bm{u}(\xi)) - \bm{R}^{\mathrm{t}}(\bm{X}_I+\bm{u}_I) - (\bm{R}_0^{\mathrm{t}}\bm{X}(\xi) - \bar{\bm{C}}) \\
&= \bm{R}^{\mathrm{t}}(\bm{X}(\xi)+\bm{u}(\xi)) - \bm{R}^{\mathrm{t}}(\bm{X}_I+\bm{u}_I) - \bm{R}_0^{\mathrm{t}}\bm{X}(\xi) + \bar{\bm{C}} \\
&= \bm{R}^{\mathrm{t}}(\bm{X}(\xi)-\bm{X}_I) 
 + \bm{R}^{\mathrm{t}}(\bm{u}(\xi)-\bm{u}_I) - \bm{R}_0^{\mathrm{t}}\bm{X}(\xi) + \bar{\bm{C}} \\
\end{aligned}
\]
Enforcing \(\bm{u}_I = \bm{0}\)
\begin{equation}
\bar{\bm{C}} = \bm{R}_0^{\mathrm{t}}\bm{X}_I
\label{eq:init-Cbar}
\end{equation}
Then we have
\begin{equation*}
\bar{\bm{u}}(\xi) = \bm{R}^{\mathrm{t}}(\bm{X}(\xi)-\bm{X}_I) 
 + \bm{R}^{\mathrm{t}}(\bm{u}(\xi)-\bm{u}_I)
 - \bm{R}_0^{\mathrm{t}}(\bm{X}(\xi) - \bm{X}_I)
\end{equation*}
% \[
% \bar{\bm{u}}_J = \bar{\bm{x}} - \bar{\bm{X}} = \bm{R}^{\mathrm{t}}\bm{u}_J - \bm{c} = \bm{R}^{\mathrm{t}}\Delta \bm{u}
% \]

\paragraph{Orientation}

We want a rotation \(\bm{R}(\varepsilon)\) (with
\(\bm{R}(0)=\bm{1}\)) that \emph{maps the original unit direction}
\(\Delta \bm{X}^{\ast} = \tfrac{\Delta \bm{X}}{\|\Delta \bm{X}\|}\) to
the \emph{new} unit direction
\(\Delta \bm{x}^{\ast}(\varepsilon) = \tfrac{\Delta \bm{x}(\varepsilon)}{\|\Delta \bm{x}(\varepsilon)\|}\).
In other words,\\
\[
\bm{R}(\varepsilon)\,\Delta \bm{X}^{\ast} 
\;=\;
\Delta \bm{x}^{\ast}(\varepsilon).
\]
That is, \(\bm{R}\) takes the original unit direction vector to
the new unit direction vector.
%
If \(\Delta \bm{X}\) and \(\Delta \bm{x}\) are \emph{not} collinear, one
can build \(\bm{R}\) via the axis--angle representation. 
Define

\[\Delta \bm{X}^\ast \triangleq \Delta \bm{X} / \|\Delta \bm{X}\|
\qquad\text{ and }\qquad
\Delta \bm{x}^\ast \triangleq \Delta \bm{x} / \|\Delta \bm{x}\|.
\]
%
The rotation axis is the unit vector 
\begin{equation}
\AlignAxis \;=\; \frac{\Delta \bm{X}^\ast \times \Delta \bm{x}^\ast}{\|\Delta \bm{X}^\ast \times \Delta \bm{x}^\ast\|}\,,
\label{eq:def-align-axis}
\end{equation}
and the rotation angle \(\theta\) is 
\begin{equation}
\theta \;=\; \arccos \bigl(\Delta \bm{X}^\ast \cdot \Delta \bm{x}^\ast\bigr)
\label{eq:arccos-dX-dx}
\end{equation}
Then Rodrigues' formula (\eqref{eq:ExpSO3}) furnishes 
\[
\bm{R} 
\;=\; 
\bm{1} 
\;+\; [\AlignAxis]_\times \,\sin \theta 
\;+\; [\AlignAxis]_\times^2 \,(1 - \cos \theta)
\]

\paragraph{Remarks}

\begin{itemize}
\tightlist
\item
  If \(\Delta \bm{x}\) is exactly parallel to \(\Delta \bm{X}\) (same
  direction), the angle is zero, and \(\bm{R}\) is the identity.\\
\item
  If they are \emph{antiparallel}, then \(\theta=\pi\), and one only
  needs an axis perpendicular to \(\Delta \bm{X}\) and \(\Delta \bm{x}\).
\end{itemize}


\subsubsection{Summary}

Define 
\begin{equation*}
    \bm{c} = \operatorname{Shift}(\chi_{I}, \chi_J)
    \qquad\text{ and }\qquad
    \bm{R} = \operatorname{Exp}\circ \operatorname{Axis}(\Delta \bm{X},\; \Delta \bm{X} + \Delta \bm{u})
\end{equation*}
where
\begin{equation*}
\begin{aligned}
\operatorname{Axis} \colon \mathbb{R}^3\times\mathbb{R}^3 &\rightarrow \mathbb{R}^3 \\
\Delta \bm{X},\,\Delta\bm{x}
&\mapsto \frac{\Delta \bm{X}^\ast \times \Delta \bm{x}^\ast}{\|\Delta \bm{X}^\ast \times \Delta \bm{x}^\ast\|} \,
\arccos \bigl(\Delta \bm{X}^\ast \cdot \Delta \bm{x}^\ast\bigr)
\end{aligned}
\end{equation*}


\begin{ambox}[Axial Transform]{box:axial-transform}

\small
\begin{enumerate}
    \item Rotation parameters
    \[
    \begin{pmatrix}
    \bm{x} \\
    \bm{\theta} 
    \end{pmatrix}
    \mapsto \begin{pmatrix}
      \bm{x} \\
      \operatorname{Exp} \bm{\theta}
    \end{pmatrix}
    \]
    \item Global offsets 
    \[
    \begin{pmatrix}
    \bm{u}(\xi_a) \\ \FrameRotation(\xi_a) 
    \end{pmatrix}
    \mapsto 
    \begin{pmatrix}
        \bm{u}(\xi_a) + \FrameRotation(\xi_a) \bm{r}_a^{}\\
        \FrameRotation(\xi_a)
    \end{pmatrix}
    \]
    \[
    \begin{pmatrix}
        \bm{v}_x\\ \bm{v}_{\Lambda}
    \end{pmatrix}
    \mapsto 
    \begin{pmatrix}
        
    \end{pmatrix}
    \]

    \item Isometry
    \[
    \bm{x}(\xi),\FrameRotation(\xi) \mapsto \begin{pmatrix}
        \bm{R}^{\mathrm{t}}(\bm{x}(\xi) -\bm{c})\\
        \bm{R}^{\mathrm{t}}\FrameRotation (\xi)
    \end{pmatrix}
    \]
    with \(\bm{c}\). Translation increments from \eqref{eq:iso-u}:
    \[
    \begin{aligned}
    \bm{u}(\xi)&\;\mapsto\; \bm{R}^{\mathrm{t}}(\bm{u}(\xi) - \Delta \bm{c}) + \Delta \bm{R}^{\mathrm{t}}\bm{X}(\xi) \\
    \bm{v}_x,\bm{v}_{\Lambda} &\mapsto \begin{pmatrix}
      \bm{R}^{\mathrm{t}}\left(\bm{v}_x - \bm{\psi}_x \right) \\
    \end{pmatrix}
    \end{aligned}
    \]

    \item Local logarithm
    \[
      \FrameRotation \mapsto \bm{\theta} \triangleq \operatorname{Log} \FrameRotation
    \]
    \[
    \bm{v}_{\Lambda} \mapsto \operatorname{dLog}(\bm{\theta}) \bm{v}_{\Lambda}
    \qquad 
    \bm{m} \mapsto \operatorname{dLog}(\bm{\theta})^{\mathrm{t}} \bm{m}
    \]
\end{enumerate}
\end{ambox}

\subsection{Linearization}

Consider small displacements of the endpoints, proportional to a small parameter \(\varepsilon\). 
Specifically, let 
\[
\bm{u}_I(\varepsilon) \;=\; \varepsilon\,\bm{u}_I, 
\quad
\bm{u}_J(\varepsilon) \;=\; \varepsilon\,\bm{u}_J,
\]
so that the new line endpoints become 
\[
\bm{X}_I + \bm{u}_I(\varepsilon)
\quad\text{and}\quad
\bm{X}_J + \bm{u}_J(\varepsilon).
\] 
Then the new direction vector is 
\begin{equation}
    \Delta \bm{x}(\varepsilon)
    \;=\;
    \bigl(\bm{X}_J + \varepsilon\,\bm{u}_J\bigr)
    \;-\;
    \bigl(\bm{X}_I + \varepsilon\,\bm{u}_I\bigr)
    \;=\;
    \Delta \bm{X} \;+\; \varepsilon\,\Delta \bm{u},
    \quad\text{ with }\quad
    \Delta \bm{u} \triangleq \bm{u}_J - \bm{u}_I.
\label{eq:order-Dx}
\end{equation} 

\subsubsection{Translation}

\subsubsection{Orientation}
At \(\varepsilon=0\), clearly \(\Delta \bm{x}(0) = \Delta \bm{X}\). 
The \emph{unit} direction under the motion is 
\[
\Delta \bm{x}^*(\varepsilon)
\; \triangleq \;
\frac{\Delta \bm{x}(\varepsilon)}{\|\Delta \bm{x}(\varepsilon)\|}.
\]
Recall the \emph{rotation angle} \(\theta(\varepsilon)\) and \emph{rotation axis} \(\AlignAxis(\varepsilon)\) are defined via \eqref{eq:arccos-dX-dx} and \eqref{eq:def-align-axis}, respectively:
  \[
    \theta(\varepsilon)
    \;=\;
    \arccos~\Bigl(
      \Delta \bm{X}^*\cdot \Delta \bm{x}^*(\varepsilon)
    \Bigr),
    \qquad
    \AlignAxis(\varepsilon)
    \;=\;
    \frac{
      \Delta \bm{X}^*\times \Delta \bm{x}^*(\varepsilon)
    }{
      \bigl\|\Delta \bm{X}^*\times \Delta \bm{x}^*(\varepsilon)\bigr\|
    }.
\] 
At \(\varepsilon=0\), these expressions are not valid because
\(\Delta \bm{x}^*(0) \equiv \Delta \bm{X}^*\).  
The exponential (\eqref{eq:ExpSO3}) then furnishes for \(\varepsilon\neq 0\): 
\[
\bm{R}(\varepsilon)
\;=\;
\bm{1}
\;+\;
\sin \theta(\varepsilon) \,
[\,\AlignAxis(\varepsilon)\,]_{\times}
\;+\;
\bigl(1-\cos\theta(\varepsilon)\bigr)\,
[\,\AlignAxis(\varepsilon)\,]_{\times}^2.
\] 
This \(\bm{R}(\varepsilon)\) is well‐defined for each
\(\varepsilon\neq 0\). 
At \(\varepsilon=0\), \(\Delta \bm{x}^*(0)=\Delta \bm{X}^*\),
so \(\theta(0) = 0\) and
\(\|\Delta \bm{X}^*\times \Delta \bm{x}^*(0)\|=0\). 
In that sense, \(\AlignAxis(0)\) is not defined. Nevertheless, \(\bm{R}(0)=\bm{1}\).
%
We now want the derivative 
\[
\left.\frac{\mathrm{d}}{\mathrm{d}\varepsilon}\,\bm{R}(\varepsilon)\right|_{\varepsilon=0}.
\]
%

\paragraph{Unit Axis}

We begin by obtaining the first order approximation of the unit
direction. The norm of \(\Delta \bm{x}(\varepsilon)\) is \[
  \|\Delta \bm{x}(\varepsilon)\|
  \;=\;
  \|\Delta \bm{X} + \varepsilon\,\Delta \bm{u}\|
\] 
For small \(\varepsilon\), expand to first order:
\[
  \|\Delta \bm{X} + \varepsilon\,\Delta \bm{u}\|
  \;=\;
  \|\Delta \bm{X}\|\;
  \sqrt{\,1 
           \;+\; 
           2\,\varepsilon\,\frac{\Delta \bm{X}}{\|\Delta \bm{X}\|^2}\cdot \Delta \bm{u}
           \;+\; O(\varepsilon^2)
         }.
\] 
Hence, to first order: 
\[
  \|\Delta \bm{x}(\varepsilon)\| 
  \;=\;
  \|\Delta \bm{X}\|
  \;+\;
  \varepsilon\,
  \frac{\Delta \bm{X}}{\|\Delta \bm{X}\|}\cdot \Delta \bm{u}
  \;+\; 
  O(\varepsilon^2).
\]
%
Therefore, 
\[
  \Delta \bm{x}^\ast(\varepsilon)
  \;=\;
  \frac{\Delta \bm{x}(\varepsilon)}{\|\Delta \bm{x}(\varepsilon)\|}
  \;=\;
  \frac{\Delta \bm{X} + \varepsilon\,\Delta \bm{u}}
       {\|\Delta \bm{X}\| 
        + \varepsilon\,\frac{\Delta \bm{X}}{\|\Delta \bm{X}\|}\cdot \Delta \bm{u}
       }
  \;=\;
  \bigl(\Delta \bm{X}^\ast + \varepsilon\,\tfrac{\Delta \bm{u}}{\|\Delta \bm{X}\|}\bigr)\;\left[\,
     1 
     \;-\;
     \varepsilon\,\tfrac{\Delta \bm{X}^\ast}{\|\Delta \bm{X}\|}\cdot \Delta \bm{u}
     \;+\;
     O(\varepsilon^2)
  \right].
\] 
Carrying out the product to first order in \(\varepsilon\):

\[
  \Delta \bm{x}^\ast(\varepsilon)
  \;=\;
  \Delta \bm{X}^\ast
  \;+\;
  \varepsilon
  \biggl(
     \tfrac{\Delta \bm{u}}{\|\Delta \bm{X}\|}
     \;-\;
     (\Delta \bm{X}^\ast \!\cdot \Delta \bm{u})\;\Delta \bm{X}^\ast \;\frac{1}{\|\Delta \bm{X}\|}
  \biggr)
  \;+\;
  O(\varepsilon^2).
\] 
Collecting terms,
\[
  \Delta \bm{x}^\ast(\varepsilon)
  \;=\;
  \Delta \bm{X}^\ast
  \;+\;
  \varepsilon\,
  \frac{1}{\|\Delta \bm{X}\|}\bigl[\Delta \bm{u} 
           \;-\; 
           (\Delta \bm{X}^\ast\!\cdot \Delta \bm{u})\,\Delta \bm{X}^\ast
    \bigr]
  \;+\;
  O(\varepsilon^2).
\]
and
\[
\Delta \bm{x}^\ast(\varepsilon)
  \;=\; ?\left(
  \bm{1}
  \;+\;
  \varepsilon\,
  \frac{1}{\|\Delta \bm{X}\|}\bigl[
      \Delta \bm{u} \otimes \Delta\bm{X}^*
      \;-\; 
      (\Delta \bm{X}^\ast\!\cdot \Delta \bm{u})\,\bm{1} %\Delta \bm{X}^\ast
    \bigr]
  \right) \Delta \bm{X}^\ast
  \;+\;
  O(\varepsilon^2).
\]
Thus the \emph{change} in the unit direction to first order is 
\begin{equation}
  \Delta \mathbf{e}(\varepsilon)
  \;\triangleq\;
  \Delta \bm{x}^\ast(\varepsilon) \;-\; \Delta \bm{X}^\ast
  \;=\;
  \varepsilon\,
  \frac{
    \Delta \bm{u} 
           - (\Delta \bm{X}^\ast\!\cdot \Delta \bm{u})\,\Delta \bm{X}^\ast
  }{\|\Delta \bm{X}\|}
  \;+\;
  O(\varepsilon^2).
  \label{eq:delta-basis}
\end{equation}

\begin{marsdenbox}
We have 
\[
  \Delta \bm{x}(\varepsilon)
  \;=\;
  \Delta \bm{X}
  \;+\;
  \varepsilon\,\Delta \bm{u},
  \quad
  \|\Delta \bm{x}(\varepsilon)\|
  \;=\;
  \|\Delta \bm{X} + \varepsilon\,\Delta \bm{u}\|.
\] 
Hence \[
  \Delta \bm{x}^*(\varepsilon)
  \;=\;
  \frac{\Delta \bm{X} + \varepsilon\,\Delta \bm{u}}{\|\Delta \bm{X} + \varepsilon\,\Delta \bm{u}\|}.
\] 
Define two scalar functions \[
  g(\varepsilon) \;=\; \Delta \bm{X} + \varepsilon\,\Delta \bm{u}, 
  \qquad \text{ and }\qquad 
  h(\varepsilon) \;=\; \|g(\varepsilon)\|.
\] 
Then
\(\Delta \bm{x}^*(\varepsilon) = g(\varepsilon)/h(\varepsilon)\). 
The
derivative with respect to \(\varepsilon\) is 
\[
  \frac{\mathrm{d}}{\mathrm{d}\varepsilon}\,\frac{g(\varepsilon)}{h(\varepsilon)}
  \;=\;
  \frac{g'(\varepsilon)}{h(\varepsilon)}
  \;-\;
  \frac{g(\varepsilon)\,h'(\varepsilon)}{h(\varepsilon)^2}.
\] 
At \(\varepsilon=0\), 
\[
  g(0)=\Delta \bm{X},\quad
  g'(0)=\Delta \bm{u},
  \quad
  h(0)=\|\Delta \bm{X}\|,
  \quad
  h'(0)
  \;=\;
  \left.\frac{\mathrm{d}}{\mathrm{d}\varepsilon}\|g(\varepsilon)\|\right|_{\varepsilon=0}
  \;=\;
  \frac{\Delta \bm{X}}{\|\Delta \bm{X}\|} \cdot \Delta \bm{u}.
\] 
Thus 
\[
\left.\frac{\mathrm{d}}{\mathrm{d}\varepsilon}\Delta \bm{x}^*(\varepsilon)\right|_{\varepsilon=0}
  \;=\;
  \frac{\Delta \bm{u}}{\|\Delta \bm{X}\|}
  \;-\;
  \frac{\Delta \bm{X}}{\|\Delta \bm{X}\|}\,
  \frac{\Delta \bm{X}}{\|\Delta \bm{X}\|^2}\cdot \Delta \bm{u}
\]
and upon simplification:
\begin{equation}
  \;=\;
  \frac{1}{\|\Delta \bm{X}\|}
  \Bigl[
    \bm{1}
    \;-\;
    \Delta \bm{X}^* \otimes \Delta \bm{X}^*
  \Bigr]\Delta \bm{u}.
  \label{eq:linear-Dxstar}
\end{equation}
%
\end{marsdenbox}

Note that \eqref{eq:linear-Dxstar} implies 
\[
  \Delta \bm{X}^* \,\cdot\, \Delta \bm{x}^*(\varepsilon)
  \;=\;
  1
  \;+\;
  \varepsilon
  \underbrace{
    \frac{
      \Delta \bm{X}^*\!\cdot \Delta \bm{u}
      \;-\;
      \bigl(\Delta \bm{X}^*\!\cdot\Delta \bm{u}\bigr)\, \Delta \bm{X}^* \cdot \Delta \bm{X}^*
    }{\|\Delta \bm{X}\|}
  }_{=\,0}
  \;+\;
  O(\varepsilon^2).
\] 
where the first-order term \(\Delta \mathbf{e} \cdot \Delta \bm{X}^*\) is zero because 
\[
  \Delta \bm{X}^*\cdot \Delta \bm{u}
  \;-\;
  (\Delta \bm{X}^*\!\cdot\Delta \bm{u})\,(\Delta \bm{X}^*\!\cdot\Delta \bm{X}^*)
  \;=\;
  \Delta \bm{X}^*\!\cdot \Delta \bm{u}
  \;-\;
  \Delta \bm{X}^*\!\cdot\Delta \bm{u}
  \;=\;
  0.
\] 
Hence
\begin{equation}
  \Delta \bm{X}^* \,\cdot\, \Delta \bm{x}^*(\varepsilon)
  \;=\;
  1
  \;+\;
  O(\varepsilon^2).
  \label{eq:order-DXDx}
\end{equation}

\paragraph{\texorpdfstring{Expand \(\theta(\varepsilon)\)}{Expand theta(e)}}

Recall \eqref{eq:arccos-dX-dx}:
\[
  \theta(\varepsilon)
  \;=\;
  \arccos \bigl(
    \Delta \bm{X}^*\cdot \Delta \bm{x}^*(\varepsilon)
  \bigr),
\] 
Therefore, to first order, \eqref{eq:order-DXDx} implies
\[
  \theta(\varepsilon)
  \;=\;
  \arccos \bigl(1 + O(\varepsilon^2)\bigr)
  \;=\;
  O(\varepsilon^2).
\]
That is, \textbf{there is \emph{no} first-order change} in the dot
product---so the angle is actually \emph{second-order} small in
\(\varepsilon\). Equivalently, the \emph{component} of \(\Delta \bm{u}\)
parallel to \(\Delta \bm{X}\) does \emph{not} create a first-order
\emph{rotation}; it only contributes to length changes or second-order
tilt.
From this one can show (by direct Taylor expansion of
\(\Delta \bm{x}^*(\varepsilon)\) around \(\varepsilon=0\)) that \[
  \Delta \bm{X}^*\cdot \Delta \bm{x}^*(\varepsilon)
  \;=\;
  1 
  \;-\;
  \frac{\bigl\|\Delta \bm{X}^*\times \Delta \bm{x}^*(\varepsilon)\bigr\|^2}{2}
  \quad\text{(at small }\varepsilon\text{)},
\] and that
\(\bigl\|\Delta \bm{X}^*\times \Delta \bm{x}^*(\varepsilon)\bigr\|\) is
typically \emph{proportional} to \(\varepsilon\), \emph{unless} the
displacement \(\Delta \bm{u}\) has a purely parallel component (in which
case the line is just stretched, so \(\theta(\varepsilon)\approx 0\) to
second order, which still yields a well‐defined limit).

\begin{itemize}
\tightlist
\item
  \textbf{Case (a):} If \(\Delta \bm{u}\) is \emph{not} perfectly
  parallel to \(\Delta \bm{X}\), then \(\theta(\varepsilon)\) is
  typically \(O(\varepsilon)\).\\
\item
  \textbf{Case (b):} If \(\Delta \bm{u}\) is parallel, you get no
  first‐order rotation; \(\theta(\varepsilon)\) is \(O(\varepsilon^2)\).
\end{itemize}
%
Either way, we know that as \(\varepsilon\to0\),
\(\theta(\varepsilon)\to 0\). 
Then 
\begin{itemize}
\tightlist
\item
  \(\sin \theta(\varepsilon)  = \theta(\varepsilon) + O(\theta^3)\).\\
\item
  \(1-\cos[\theta(\varepsilon)] = \tfrac12\,\theta^2(\varepsilon) + O(\theta^4)\).
\end{itemize}
%
Thus

\[
  \frac{\sin \theta(\varepsilon) }{\varepsilon}
  \;=\;
  \frac{\theta(\varepsilon)}{\varepsilon}
  \;+\;
  O\!\bigl(\theta^3/\varepsilon\bigr),
  \quad
  \frac{1-\cos[\theta(\varepsilon)]}{\varepsilon}
  \;=\;
  \frac{\tfrac12\theta^2(\varepsilon)}{\varepsilon}
  \;+\;
  O\!\bigl(\theta^4/\varepsilon\bigr).
\] If \(\theta(\varepsilon)\) is \emph{linear} in \(\varepsilon\) at
small \(\varepsilon\) (the usual perpendicular‐displacement case), then
\(\theta(\varepsilon)/\varepsilon \to \text{constant}\), and
\(\theta^2(\varepsilon)/\varepsilon\to 0\). If \(\theta(\varepsilon)\)
is \emph{second‐order} small, then
\(\theta(\varepsilon)/\varepsilon \to0\) anyway. In \emph{either}
scenario,

\begin{enumerate}
\def\labelenumi{\arabic{enumi}.}
\tightlist
\item
  \(\sin \theta(\varepsilon) /\varepsilon\) is \emph{bounded} as
  \(\varepsilon \to 0\) and has a finite limit (which may be 0 in the
  purely parallel case).\\
\item
  \((1-\cos \theta(\varepsilon))/\varepsilon \to 0\).
\end{enumerate}

\paragraph{\texorpdfstring{Expand \(\AlignAxis(\varepsilon)\)}{Expand i(e)}}

On the other hand, 
\[
  \AlignAxis(\varepsilon)
  \;=\;
  \frac{\Delta \bm{X}^* \times \Delta \bm{x}^*(\varepsilon)}{
        \|\Delta \bm{X}^* \times \Delta \bm{x}^*(\varepsilon)\|
       }.
\]
Let's look at the cross product in the numerator: \[
  \Delta \bm{X}^* \times \Delta \bm{x}^*(\varepsilon)
  \;=\;
  \Delta \bm{X}^* \;\times\;
  \Bigl(\Delta \bm{X}^* + \varepsilon\,\Delta\Bigr)
  \;=\;
  \Delta \bm{X}^* \times \Delta \bm{X}^*
  \;+\;
  \varepsilon\,\Delta \bm{X}^*\times\Delta,
\]
where
\(\Delta = \tfrac{1}{\|\Delta \bm{X}\|}\Bigl[\Delta \bm{u} - (\Delta \bm{X}^*\cdot\Delta \bm{u})\,\Delta \bm{X}^*\Bigr]\).
But \(\Delta \bm{X}^* \times \Delta \bm{X}^* = \mathbf{0}\), so \[
  \Delta \bm{X}^* \times \Delta \bm{x}^*(\varepsilon)
  \;=\;
  \varepsilon\,\Delta \bm{X}^*\times\Delta
  \;=\;
  O(\varepsilon).
\]
Hence the magnitude of that cross product is also \(O(\varepsilon)\).
Dividing by its own magnitude, \(\AlignAxis(\varepsilon)\) remains a
\emph{well-defined} unit vector for \(\varepsilon\neq 0\), but it
basically becomes

\[
  \AlignAxis(\varepsilon)
  \;\approx\;
  \frac{\Delta \bm{X}^*\times \Delta}{\|\Delta \bm{X}^*\times \Delta\|}.
\] In fact, to \emph{first} order in \(\varepsilon\), this is just the
direction perpendicular to \(\Delta \bm{X}^*\) and \(\Delta\). 


\subsection{Linearization (3)}

By construction, the rotation \(\bm{R}(\varepsilon)\) carries the
original unit direction \(\Delta \bm{X}^*\) to the new unit direction
\(\Delta \bm{x}^*(\varepsilon)\). So \[
  \bm{R}(\varepsilon)\,\Delta \bm{X}^*
  \;=\;
  \Delta \bm{x}^*(\varepsilon).
\] Differentiate both sides w.r.t. \(\varepsilon\). At
\(\varepsilon=0\), \(\bm{R}(0)=\bm{1}\). 
We get 
\[
  \left.\frac{\mathrm{d}}{\mathrm{d}\varepsilon}\bm{R}(\varepsilon)\right|_{\varepsilon=0}\,\Delta \bm{X}^*
  \;+\;
  \bm{R}(0)\;
  \left.\frac{\mathrm{d}}{\mathrm{d}\varepsilon}\bigl[\Delta \bm{X}^*\bigr]\right|_{\varepsilon=0}
  \;=\;
  \left.\frac{\mathrm{d}}{\mathrm{d}\varepsilon}\Delta \bm{x}^*(\varepsilon)\right|_{\varepsilon=0}.
\]
But \(\Delta \bm{X}^*\) is a fixed (constant) vector, so its
derivative is zero. Also \(\bm{R}(0)=\bm{1}\). Thus \[
  \left.\frac{\mathrm{d}}{\mathrm{d}\varepsilon}\bm{R}(\varepsilon)\right|_{\varepsilon=0}\,\Delta \bm{X}^*
  \;=\;
  \left.\frac{\mathrm{d}}{\mathrm{d}\varepsilon}\Delta \bm{x}^*(\varepsilon)\right|_{\varepsilon=0}.
\] 
Using \eqref{eq:linear-Dxstar} for 
\(\displaystyle \bigl.\tfrac{d}{d\varepsilon}\Delta \bm{x}^*(\varepsilon)\bigr|_{\varepsilon=0}\) furnishes
%
\[
  \left.\frac{\mathrm{d}}{\mathrm{d}\varepsilon}\bm{R}(\varepsilon)\right|_{\varepsilon=0}\Delta \bm{X}^*
  \;=\;
  \frac{1}{\|\Delta \bm{X}\|}
  \Bigl[
    \Delta \bm{u}
    - (\Delta \bm{X}^*\!\cdot \Delta \bm{u})\,\Delta \bm{X}^*
  \Bigr].
\] Crucially, that right-hand side is \emph{orthogonal} to
\(\Delta \bm{X}^*\), which shows that
\(\bigl.\tfrac{d}{d\varepsilon}\bm{R}(\varepsilon)\bigr|_{\varepsilon=0}\)
acts like a \emph{skew-symmetric} operator on \(\Delta \bm{X}^*\). 
%
Indeed, a rank-1 check shows there must exist some vector
\(\bm{\omega}\) such that for \emph{all} \(\mathbf{x}\), \[
  \left.\frac{\mathrm{d}}{\mathrm{d}\varepsilon}\bm{R}(\varepsilon)\right|_{\varepsilon=0}\mathbf{x}
  \;=\;
  \bm{\omega}\times \mathbf{x}.
\] 
Applying it to \(\Delta \bm{X}^*\) gives 
\[
  \bm{\omega}\times \Delta \bm{X}^*
  \;=\;
  \frac{1}{\|\Delta \bm{X}\|}
  \Bigl[
    \Delta \bm{u}
    - (\Delta \bm{X}^*\!\cdot \Delta \bm{u})\,\Delta \bm{X}^*
  \Bigr].
\] 
One can check that
\(\bm{\omega} = \tfrac{1}{\|\Delta \bm{X}\|}\,\bigl(\Delta \bm{X}^*\times \Delta \bm{u}\bigr)\)
satisfies exactly that. 
Thus, finally, 
\begin{equation}
  \boxed{
  \frac{\mathrm{d}}{\mathrm{d}\varepsilon}\bm{R}(\varepsilon)\Big|_{\varepsilon=0}
  \;=\;
  \bm{\omega}^{\times},
  \qquad
  \bm{\omega}
  \;=\;
  \frac{\Delta \bm{X}\times \Delta \bm{u}}{\|\Delta \bm{X}\|^2}.
  }
\end{equation}



\begin{equation*}
\boxed{
  R'(0) = \lim_{\varepsilon\to 0}\,
  \frac{\bm{R}(\varepsilon) - \bm{1}}{\varepsilon}
  \;=\;
  \left[
    \frac{\Delta \bm{X}\times \Delta \bm{u}}{\|\Delta \bm{X}\|^2}
  \right]^{\times}.
}
\end{equation*}

% Equivalently,
% \[
%   \boxed{
%   R'(0)
%   \;=\;
%   [\bm{\omega}]_{\times},
%   \quad
%   \text{where }
%   \bm{\omega}
%   \;=\;
%   \frac{\Delta \bm{X}\times \Delta \bm{u}}{\|\Delta \bm{X}\|^2}.
%   }
% \]


\subsection{Summary}


\begin{itemize}
\tightlist
\item \(\bm{\omega}\) represents an \emph{infinitesimal rotation} around the axis
given by \(\Delta \bm{X} \times \Delta \bm{u}\). Geometrically, it is
precisely the smallest rotation needed to tip \(\Delta \bm{X}\) into
\(\Delta \bm{X} + \varepsilon \,\Delta \bm{u}\), modulo normalization.

\item
  \(\Delta \bm{X}^*\times \Delta \bm{x}^*(\varepsilon)\) is an
  \(O(\varepsilon)\) vector (unless purely parallel, in which case it
  stays 0 up to second order---but then \(\sin[\theta]\approx 0\) at
  first order, so the product still can have a finite limit).\\
\end{itemize}
%
Stated otherwise, if \(\bm{u}_{\perp}\) is the component of \(\Delta \bm{u}\)
perpendicular to \(\Delta \bm{X}^*\), then for small \(\varepsilon\),

\[
  \Delta \bm{X}^*\times \Delta \bm{x}^*(\varepsilon)
  \;\approx\;
  \varepsilon\,\Bigl[\Delta \bm{X}^*\times \bm{u}_{\perp}\Bigr]
  \;+\;
  O(\varepsilon^2).
\] Hence \[
  \AlignAxis(\varepsilon)
  \;=\;
  \frac{\Delta \bm{X}^*\times \Delta \bm{x}^*(\varepsilon)}{
        \|\Delta \bm{X}^*\times \Delta \bm{x}^*(\varepsilon)\|}
  \;\to\;
  \frac{\Delta \bm{X}^*\times \bm{u}_{\perp}}{\|\Delta \bm{X}^*\times \bm{u}_{\perp}\|},
  \quad\text{as }\varepsilon\to 0,
\] if \(\bm{u}_{\perp}\neq 0\). (If \(\bm{u}_{\perp}=0\), the entire
rotation is second‐order small and the limit for
\(\sin \theta(\varepsilon) /\varepsilon\cdot\AlignAxis(\varepsilon)\)
still goes to 0.)


%
The essence is that at small angles, 
\(\sin\theta(\varepsilon)\approx\|\Delta \bm{X}^*\times \Delta \bm{x}^*\|\approx O(\varepsilon)\),\\
while \(\,1-\cos\theta(\varepsilon)\approx O(\varepsilon^2)\), and the
rotation axis becomes the direction of
\(\Delta \bm{X}^*\times(\Delta \bm{X}^*+\dots)\). 
Collecting these expansions yields
\[
\bm{\omega}=\frac{\Delta \bm{X}\times \Delta \bm{u}}{\|\Delta \bm{X}\|^2}
\].

\vspace{5ex}

\begin{ambox}[Linear Transform]{box:linear-transform}
\begin{enumerate}
    \item Rotation Parameters
    \item Global offset
    \[
    \bm{u}_a,\bm{w}_a \mapsto \begin{pmatrix}
        \bm{u}_a - \bm{r}_a^{\mathrm{g}}\times \bm{w}_a \\
        \bm{w}_a
    \end{pmatrix}
    \qquad\qquad
    \bm{n}_a,\bm{m}_a \mapsto \begin{pmatrix}
        \bm{n}_a \\
        \bm{m}_a + \bm{r}_a \times\bm{n}_a
    \end{pmatrix}
    \]
    \item 
    \[
    \bm{u}_a,\bm{w}_a \mapsto \begin{pmatrix}
        \bm{R}^{\mathrm{t}}\bm{u}_a \\
        \bm{R}^{\mathrm{t}}\bm{w}_a
    \end{pmatrix}
    \qquad\qquad
    \bm{n}_a,\bm{m}_a \mapsto \begin{pmatrix}
        \bm{R}\bm{n}_a \\
        \bm{R}\bm{m}_a
    \end{pmatrix}
    \]
    \item Local offset
    \item Isometry
    \[
    \bm{u}_a,\FrameRotation \mapsto \begin{pmatrix}
        \bm{u}_{a} + \bm{c} + \Delta\bm{x}_a \times \bar{\bm{\omega}} \\
        \bm{R}^{\mathrm{t}}\FrameRotation
    \end{pmatrix}
    \]
    \[
    \bm{v}_{xa},\bm{v}_{\Lambda a} \mapsto \begin{pmatrix}
        \bm{v}_{xa} + \bm{c} + \Delta\bm{x}_{a} \times \bar{\bm{\omega}} \\
        \bm{v}_{\Lambda a} - \bar{\bm{\omega}}
    \end{pmatrix}
    \qquad\qquad
    \bm{n}_a,\bm{m}_a \mapsto \begin{pmatrix}
        \bm{n}_a - \frac{\delta}{L}\FrameBasis \times \bar{\bm{m}} \\
        \bm{m}_a
    \end{pmatrix}
    \]
    with \(\Delta\bm{u}\), \(\bm{c} \triangleq - \bm{u}_I\) and \(\bar{\bm{\omega}} \triangleq \frac{1}{L} \FrameBasis \times \Delta \bm{u}\)

    \item Local Logarithm

\end{enumerate}
\end{ambox}

\vspace{5ex}
\begin{ambox}[Linear Transform]{box:linear-transform}
\begin{enumerate}
    \item[0] Node logarithm
    \[
    \bm{u}_a,\FrameRotation_a \mapsto \begin{pmatrix}
        \bm{u}_a \\
        \operatorname{Log}\FrameRotation_a
    \end{pmatrix}
    \qquad\qquad
    \bm{n}_a,\bm{m}_a \mapsto \begin{pmatrix}
        \bm{n}_a \\
        \bm{m}_a + \bm{r}_a \times\bm{n}_a
    \end{pmatrix}
    \]
    \item Global offset
    \[
    \bm{u}_a,\bm{w}_a \mapsto \begin{pmatrix}
        \bm{u}_a - \bm{r}_a^{\mathrm{g}}\times \bm{w}_a \\
        \bm{w}_a
    \end{pmatrix}
    \qquad\qquad
    \bm{n}_a,\bm{m}_a \mapsto \begin{pmatrix}
        \bm{n}_a \\
        \bm{m}_a + \bm{r}_a \times\bm{n}_a
    \end{pmatrix}
    \]
    \item 
    \[
    \bm{u}_a,\bm{w}_a \mapsto \begin{pmatrix}
        \bm{R}^{\mathrm{t}}\bm{u}_a \\
        \bm{R}^{\mathrm{t}}\bm{w}_a
    \end{pmatrix}
    \qquad\qquad
    \bm{n}_a,\bm{m}_a \mapsto \begin{pmatrix}
        \bm{R}\bm{n}_a \\
        \bm{R}\bm{m}_a
    \end{pmatrix}
    \]
    \item Local offset
    \item 
    \[
    \bm{u}_a,\bm{w}_a \mapsto \begin{pmatrix}
        (\bm{u}_{a \in [I,J]} - \bm{u}_I) 
        - \bar{\bm{\omega}}\times (\Delta\bm{x}_a + \bm{u}_{a \notin [I,J]}) \\
        \bm{w}_a - \bar{\bm{\omega}}
    \end{pmatrix}
    \]
    with \(\Delta\bm{u}\), and \(\bar{\bm{\omega}} \triangleq \frac{1}{L} \FrameBasis \times \Delta \bm{u}\)
\end{enumerate}
\end{ambox}



\subsection{Axial System C2}
This formulation comes from \cite{crisfield1990consistent}. 
A ``mean rotation matrix" $\bm{R}_{\mathrm{aux}}$, is first formed as follows

$$
\begin{aligned}
\bm{R}_{\mathrm{aux}}
&=\Delta R\left(\frac{\omega_J-\omega_I}{2}\right) \bm{\Lambda}_I(\omega_I) \\
&=\Delta R\left(\frac{\omega_I-\omega_J}{2}\right) \bm{\Lambda}_J(\omega_J) \\
& \simeq \operatorname{Exp}\left(\frac{\omega_I+\omega_J}{2}\right)
\end{aligned}
$$

and

$$
\Delta \bm{R}(\omega_J-\omega_I)=\bm{\Lambda}_J \bm{\Lambda}_I^{\mathrm{t}}
$$


$$
(\texttt{C2})\left\{\begin{aligned}
\bar{\mathbf{d}}_1 &= \frac{\bm{x}_J-\bm{x}_I}{\left\|\bm{x}_J-\bm{x}_I\right\|} \\
\bar{\mathbf{d}}_2 &=  \operatorname{Exp}(\bm{\rho}) \mathbf{t}_2 \\
\bar{\mathbf{d}}_3 &=  \operatorname{Exp}(\bm{\rho}) \mathbf{t}_3 \\
\end{aligned}\right.
, \quad \bm{\rho} =\cos^{-1}\left(\mathbf{t}_1 \cdot \bar{\mathbf{d}}_1\right) \frac{\mathbf{t}_1 \times \bar{\mathbf{d}}_1}{\left\|\mathbf{t}_1 \times \bar{\mathbf{d}}_1\right\|}
$$

$$
\boxed{
\bm{R} \triangleq \operatorname{Exp}(\operatorname{Axis}(\bm{R}_{\mathrm{aux}}\mathbf{D}_1,\, \mathbf{D}_1)) \, \bm{R}_{\mathrm{aux}}
}
$$

\subsection{Axial System C1}

In order to facilitate taking variations, Crisfield suggests approximating the 
correcting rotation $\operatorname{Exp}(\bm{\rho})$ by a ``midpoint" formula as follows:

$$
(\texttt{C1})\left\{\begin{aligned}
\bar{\mathbf{d}}_1 &= \frac{\bm{x}_J-\bm{x}_I}{\left\|\bm{x}_J-\bm{x}_I\right\|} \\
\bar{\mathbf{d}}_2 &= \mathbf{t}_2 - \frac{1}{2} \mathbf{t}_2 \cdot \bar{\mathbf{d}}_1  \left(\bar{\mathbf{d}}_1+\mathbf{t}_1\right), \\
\bar{\mathbf{d}}_3 &= \mathbf{t}_3 - \frac{1}{2} \mathbf{t}_3 \cdot \bar{\mathbf{d}}_1  \left(\bar{\mathbf{d}}_1+\mathbf{t}_1\right) .
\end{aligned}\right.
$$


\subsection{Axial System B}


$$
\mathbf{t}=\frac{1}{2}\left(\mathbf{t}_1+\mathbf{t}_2\right), \quad \mathbf{t}_i=\bm{\Lambda}_i \bm{R}_0 \mathbf{D}_2, \quad i=1,2
$$

$$
\begin{cases}
\bar{\mathbf{d}}_{1} & =\dfrac{\bm{x}_J-\bm{x}_I}{\left\|\bm{x}_J-\bm{x}_I\right\|} \\
\bar{\mathbf{d}}_{3} & =\dfrac{\bar{\mathbf{d}}_1 \times \mathbf{t}}{\left\|\bar{\mathbf{d}}_1 \times \mathbf{t}\right\|}  \\
\bar{\mathbf{d}}_{2} & = \bar{\mathbf{d}}_{3} \times  \bar{\mathbf{d}}_1
\end{cases}
$$

$$
\operatorname{Align}(\FrameBasis \mapsto \frameBasis, \, D_3 \mapsto t)
= \operatorname{Exp}(\frameBasis\, \operatorname{Angle}(...)\, )
\operatorname{Exp}\bigl(\operatorname{Axis}(\FrameBasis, \frameBasis) \bigr)
$$



\subsection{Sensitivity}

\begin{equation}
\begin{aligned}
L^{\prime} & =\frac{1}{L}\left(\bm{X}_j-\bm{X}_i\right) \cdot\left(\bm{X}_2^{\prime}-\bm{X}_i^{\prime}\right) \\
\mathbf{e}_1^{\prime} & =\frac{1}{L_1^2}\left(\Delta \bm{X}^{\prime} L-\Delta \bm{X} L^{\prime}\right) \\
\mathbf{e}_2^{\prime} & =\frac{1}{\left(\left\|\mathbf{v}_2\right\|^{\prime}\right)^2} \left(\mathbf{v}_2^{\prime} \left\|\mathbf{v}_2\right\|-\mathbf{v}_2\left\|\mathbf{v}_2\right\|^{\prime}\right) \\
%
\mathbf{v}_2^{\prime} & =\mathbf{v}^{\prime} \times \mathbf{e}_1+\mathbf{v} \times \mathbf{e}_1^{\prime} \\
\left\|\mathbf{v}_2\right\|^{\prime} & =\frac{1}{\left\|\mathbf{v}_2\right\|}  \mathbf{v}_2 \cdot \mathbf{v}_2^{\prime}
\end{aligned}
\end{equation}
